\section{Метрики оценки эффективности системы}

Для анализа использовались следующие группы метрик (см. подробнее в приложении А):
\begin{itemize}
	\item \textbf{Эффективность:} среднее социальное благосостояние $\overline{SW}$ и число завершённых задач;
	\item \textbf{Пропускная способность:} средний процент принятых задач;
	\item \textbf{Временные характеристики:} средняя задержка выполнения задачи;
	\item \textbf{Справедливость:} коэффициент Джини платежей и индекс справедливости распределения;
	\item \textbf{Сходимость обучения:} среднее вознаграждение и TD-ошибка.
\end{itemize}

Интерпретация метрик следующая: большее значение $\overline{SW}$ и acceptance rate является предпочтительным, тогда как для задержки и коэффициента Джини лучшими считаются меньшие значения.
